\documentclass[pdf,blends,slideColor,colorBG]{prosper}

\usepackage{moreverb}
\usepackage{macros}

\begin{document}

\title{Le mod�le relationnel}
\author{Philippe Rigaux}
\email{rigaux@lri.fr}
\institution{LRI}

%------------------------------------------------------------
\part{Le mod�le relationnel}
%------------------------------------------------------------


\begin{tsp}{Les concepts}
{
Relation, instance et sch�ma, domaine et attributs, cl�, tuple
et base de donn�es
\begin{center}
\begin{tabular}{l|c|l}
 titre &  ann�e & genre\\
\hline
Alien & 1979 & Science-Fiction \\
Vertigo  & 1958 & Suspense \\
Volte-face & 1997 & Thriller \\
Pulp Fiction & 1995 & Policier \\
\hline
\end{tabular}
\end{center}

}\end{tsp}
%%%%%%%%%%%%%%%%%%%%%%%%%%%%%%%%%%%%%%%%
\begin{tsp}{Notre sch�ma}{

\figWithSize{officiel}{9cm}

}\end{tsp}
%%%%%%%%%%%%%%%%%%%%%%%%%%%%%%%%%%%%%%%%%%%%%%%%%%%
\begin{tsp}{La table \tble{Film}}{

\begin{center}
\begin{small}
\begin{tabular}{l|l|c|l|c|l}
\cle{idFilm} & titre & ann�e  & genre & \fk{idMES} & \fk{codePays}\\
\hline
100 & Alien & 1979 & Science Fiction &  1 & US \\
101 & Vertigo  & 1958 & Suspense & 2 & US \\
102 & Psychose & 1960 &  Suspense & 2 & US\\
103 & Kagemusha & 1980 & Drame & 3 & JP\\
104 & Volte-face & 1997 & Action & 4 & US\\
105 & Van Gogh & 1991 & Drame  & 8 & FR\\
106 & Titanic & 1997 & Drame & 6 & US\\
107 & Sacrifice & 1986 & Drame & 7 & FR\\
\hline
\end{tabular}
\end{small}
\end{center}

}\end{tsp}
%%%%%%%%%%%%%%%%%%%%%%%%%%%%%%%%%%%%%%%%%%%%%%%%%%%
\begin{tsp}{La table \tble{Artiste}}{

\begin{center}
\begin{small}
\begin{tabular}{l|l|l|c}
{\bf idArtiste} &  nom & pr�nom &  ann�eNaiss \\
\hline
1 & Scott & Ridley & 1943\\
2 & Hitchcock & Alfred & 1899\\
3 & Kurosawa & Akira & 1910\\
4 & Woo & John & 1946\\
6 & Cameron & James & 1954\\
7 & Tarkovski & Andrei & 1932\\
8 & Pialat & Maurice & 1925 \\
\hline
\end{tabular}
\end{small}
\end{center}

}\end{tsp}
%%%%%%%%%%%%%%%%%%%%%%%%%%%%%%%%%%%%%%%%
%%%%%%%%%%%%%%%%%%%%%%%%%%%%%%%%%%%%%%%%%%%%%%%%%%%
\begin{tsp}{Repr�sentation des r�les}{

\begin{center}
\begin{small}
\begin{tabular}{l|l|c|l|c|c}
\cle{idFilm} & titre & ann�e  & genre & \fk{idMES} & \fk{codePays}\\
\hline
20 & Impitoyable & 1992 & Western & 100 & USA \\
21 & Ennemi d'�tat & 1998 & Action & 102 & USA\\
\hline
\end{tabular}
\end{small}
\end{center}

\begin{center}
\begin{small}
\begin{tabular}{l|l|l|c}
\cle{idArtiste} &  nom & pr�nom &  ann�eNaiss \\
\hline
100 & Eastwood & Clint  & 1930 \\
101 & Hackman & Gene  & 1930 \\
102 & Scott & Tony  & 1930 \\
103 & Smith  & Will  & 1968 \\
\hline
\end{tabular}
\end{small}
\end{center}

}\end{tsp}
%%%%%%%%%%%%%%%%%%%%%%%%%%%%%%%%%%%%%%%%%%%%%%%%%%%
\begin{tsp}{La table \tble{R�le}}{

NB\,: notez le respect de l'association plusieurs � plusieurs.

\begin{center}
\begin{small}
\begin{tabular}{l|c|l}
\cle{\fk{idFilm}} & \cle{\fk{idActeur}}  & r�le \\
\hline
20  & 100  & William Munny\\
20 & 101  & Little Bill\\
21 & 101  & Bril\\
21 & 103  & Robert Dean\\
\hline
\end{tabular}
\end{small}
\end{center}

}\end{tsp}
%%%%%%%%%%%%%%%%%%%%%%%%%%%%%%%%%%%%%%%%%%%%%%%%%%%
\begin{tsp}{Associations ternaires}{

NB\,: la salle est une entit� {\em faible}. Notez la cl� de la relation.

\begin{center}
\begin{small}
\begin{tabular}{l|l|c}
\cle{\fk{idCin�ma}} & \cle{no} &  capacit�  \\ \hline
Le Rex & 1 & 200  \\\hline 
Kino & 1 & 130  \\\hline 
Le Rex & 2 & 180  \\\hline 
\end{tabular}
\end{small}
\end{center}

}\end{tsp}
%%%%%%%%%%%%%%%%%%%%%%%%%%%%%%%%%%%%%%%%%%%%%%%%%%%
\begin{tsp}{Associations ternaires (suite)}{

\begin{center}
\begin{small}
\begin{tabular}{c|c|c}
\cle{idHoraire} & heureD�but & heureFin \\\hline
1 & 14 & 16  \\\hline 
2 & 16 & 18  \\\hline
\end{tabular}
\end{small}
\end{center}

\begin{center}
\begin{small}
\begin{tabular}{c|c|c|c|c}
\cle{\fk{idFilm}} & \cle{\fk{nomCin�ma}} & \cle{\fk{noSalle}} & 
       \cle{\fk{idHoraire}} & tarif \\\hline
103 & Le Rex & 2 & 1 &  23  \\\hline 
103 & Le Rex & 2 & 2 & 56  \\\hline
100 & Kino & 1 & 2 & 45  \\\hline
102 & Le Rex &  2 & 1 & 50  \\\hline
\end{tabular}
\end{small}
\end{center}


}\end{tsp}
%%%%%%%%%%%%%%%%%%%%%%%%%%%%%%%%%%%%%%%%%%%%%%%%%%%
\begin{tsp}{Cr�ation de tables}
{
{\tt
\begin{small}
\begin{tabbing}
\= CREATE TABLE \= Internaute (\=email VARCHAR (50) NOT NULL,\\
\> \> \>                      nom   VARCHAR (20) NOT NULL,\\
\> \> \>                      prenom VARCHAR (20),\\
\> \> \>                      motDePasse VARCHAR (60) NOT NULL,\\
\> \> \>                      anneeNaiss INTEGER)
\end{tabbing}
\end{small}
}
Attention aux valeurs \texttt{NULL} ou \texttt{NOT NULL}

}\end{tsp}
%%%%%%%%%%%%%%%%%%%%%%%%%%%%%%%%%%%%%%%%%%%%%%%%%%%
\begin{tsp}{Avec cl� primaire}
{

{\tt
\begin{small}
\begin{tabbing}
\= CREATE TABLE \= Internaute \=\\
 \> \> (email VARCHAR (50) NOT NULL,\\
 \> \>                      nom   VARCHAR (20) NOT NULL,\\
 \> \>                      prenom VARCHAR (20),\\
 \> \>                      motDePasse VARCHAR (60) NOT NULL,\\
 \> \>                      anneeNaiss INTEGER,\\
 \> \>                      PRIMARY KEY (email))
\end{tabbing}
\end{small}
}

NB\,: la cl� peut �tre constitu�e de plusieurs attributs

}\end{tsp}
%%%%%%%%%%%%%%%%%%%%%%%%%%%%%%%%%%%%%%%%%%%%%%%%%%%
\begin{tsp}{Avec cl� primaire et clause \texttt{UNIQUE}}
{

\sqltxt{
\begin{small}
\begin{tabbing}
\= CREATE TABLE \= Artiste(\=id INTEGER  NOT NULL,\\
 \> \>                       nom VARCHAR (30) NOT NULL,\\
 \> \>                       prenom VARCHAR (30) NOT NULL,\\
 \> \>                        anneeNaiss INTEGER,\\
 \> \>                      PRIMARY KEY (id),\\
 \> \>                      UNIQUE (nom, prenom));
\end{tabbing}
\end{small}
}

Plus facile de remettre en cause une cl� secondaire\,!
}\end{tsp}

\begin{tsp}{Avec cl�s �trang�res}
{
\sqltxt{
\begin{small}
\begin{tabbing}
\= CREATE TABLE Film \=  (idFilm INTEGER NOT NULL,\\
                    \> titre    VARCHAR (50) NOT NULL,\\
                   \>  annee     INTEGER NOT NULL,\\
                    \> idMES    INTEGER,\\
                    \> codePays   INTEGER,\\
                    \> PRIMARY KEY (idFilm),\\
              \> FOREIGN KEY (idMES) REFERENCES Artiste,\\
	   \>   FOREIGN KEY (codePays) REFERENCES Pays);
\end{tabbing}
\end{small}
}
Important\,: les tables r�f�renc�es doivent �tre cr��es {\em avant}

}\end{tsp}

\end{document}
